\section{Equation, admissibility, and notation}\label{sec:framework}

The spatial torus \(\T^3=\R^3/\mathbb Z^3\) has volume one, and
\(\Omega=\T^3_x\times\R^3_v\). We write \(x=(y,z)\), with two transverse
coordinates \(y\) and one axial coordinate \(z\), and
\(v=(\bar v,v_3)\). In a periodic lift these coordinates range in
\(\R^2\times\R\). All spatially rescaled limits below are taken in the lift.
We write \(D_t=\partial_t+v\cdot\nabla_x\) for free transport.

For incoming velocities \(v,v_*\) and \(\omega\in S^2\), set
\[
a=(v-v_*)\cdot\omega,\qquad
v'=v-a\omega,\qquad v_*'=v_*+a\omega.
\]
The surface measure on \(S^2\) is its ordinary area measure. We use
\begin{equation}\label{eq:collision}
\begin{split}
Q(F,F)&=Q^+(F,F)-Q^-(F,F),\\
Q^+(F,F)(v)&=\int_{\R^3\times S^2}|a|F(v')F(v_*')\dd v_*\dd\omega,\\
Q^-(F,F)(v)&=F(v)\nu_F(v).
\end{split}
\end{equation}
Arguments \(t,x\), identical in every factor, are suppressed.
The associated symmetric gain is denoted \(P(f,g)\), so that
\(P(F,F)=Q^+(F,F)\). Its symmetry in the hard-sphere normalization is
included in Lemma~\ref{lem:carleman}. The loss frequency is
\begin{equation}\label{eq:frequency}
\nu_f(v)=2\pi\int_{\R^3}|v-w|f(w)\dd w.
\end{equation}
A superscript \(b\) on \(P,\nu,Q\) means that the same factor
\(\ind_{\{|v-v_*|<b\}}\) is inserted into gain and loss.
This cutoff is invariant under elastic collision.

Put
\[
M(v)=(2\pi)^{-3/2}e^{-|v|^2/2},\qquad G=F/M,
\]
and define
\[
\mathcal H_M(F)=\int_\Omega
 \bigl(F\log(F/M)-F+M\bigr)\dd x\dd v.
\]
The entropy dissipation is
\[
\mathcal D(F)=\frac14
\int_{\T^3\times\R^3\times\R^3\times S^2}
|a|\,(F'F_*'-FF_*)\log\frac{F'F_*'}{FF_*}
 \dd x\dd v\dd v_*\dd\omega.
\]
The nonnegative integrand is assigned its lower-semicontinuous
extended value when a factor vanishes. We use \(0|\log0|=0\).

\begin{definition}[Entropy-class solution]\label{def:admissible}
For a nonnegative datum \(F_0\) satisfying
\[
\int_\Omega F_0(1+|v|^2+|\log F_0|)\dd x\dd v<\infty,
\]
an admissible global solution is a nonnegative \(F\) with the following
properties.
\begin{enumerate}[label=\textup{(\roman*)}]
\item \(F\in C_w([0,\infty);L^1(\Omega))\), with \(F(0)=F_0\), where
weak continuity means continuity against every bounded measurable test.
On each bounded time interval,
\(\int_\Omega F(1+|v|^2+|\log F|)\) is uniformly bounded.
\item \(Q^\pm(F,F)/(1+G)\) belong to
\(L^1([0,T]\times\Omega)\) for every finite \(T\).
For every \(\beta\in C^1([0,\infty))\) with \(\beta(0)=0\) and
\(|\beta'(z)|\le C_\beta/(1+z)\), and every smooth test
\(\varphi(t,x,v)\) periodic in \(x\) and compactly supported in
\([0,\infty)\times\R^3_v\),
\[
\begin{split}
&\int_0^\infty\!\!\int_\Omega
 \bigl[M\beta(G)(\partial_t\varphi+v\cdot\nabla_x\varphi)
      +\beta'(G)Q(F,F)\varphi\bigr]\dd x\dd v\dd t\\
&\hspace{2cm}
+\int_\Omega M\beta(F_0/M)\varphi(0,x,v)\dd x\dd v=0 .
\end{split}
\]
\item With \(\rho=\int F\dd v\) and \(j=\int vF\dd v\),
\(\partial_t\rho+\nabla_x\cdot j=0\) distributionally with the initial
density \(\int F_0\dd v\). Total momentum is conserved. For every \(t\ge0\),
\[
\int_\Omega |v|^2F(t)\le\int_\Omega |v|^2F_0,\qquad
\mathcal H_M(F(t))+\int_0^t\mathcal D(F(s))\dd s
 \le\mathcal H_M(F_0).
\]
\end{enumerate}
\end{definition}

The same definition on a finite interval uses tests and traces on that
interval. We use the global existence theory
\cite{DiPernaLions1989,DiPernaLions1991,LevermoreMasmoudi2010} only to
continue solutions already constructed on a short interval. In particular,
we do not require local momentum or local energy conservation as additional
admissibility conditions.

\subsection{A velocity weight and colored components}

Fix \(1<q<2\) and \(m=100\). The weight used throughout is
\begin{equation}\label{eq:weight}
H(v)=e^{-|v|^2}\langle v\rangle^{-m}
       \bigl(1+|\bar v|^{-q}\bigr),\qquad
\langle v\rangle=(1+|v|^2)^{1/2}.
\end{equation}
Its values on the velocity axis can be assigned arbitrarily; all weighted
supremum estimates are essential estimates. It is integrable with every
polynomial velocity weight and lies in \(L^p\) for some \(p>1\).
Write \(\|f\|_H=\esssup_{x,v}|f(x,v)|/H(v)\).

The cold components are \(c_i\), their sum is \(c\), and \(h\) denotes
the hot component. Particles retain a cold label only in collisions
within that label. All other collisions create hot particles. With
\(F=c+h\), this gives
\begin{equation}\label{eq:colored}
\begin{split}
(\D)c_i&=P(c_i,c_i)-c_i\nu_F,\\
(\D)h&=2P(c,h)+P(h,h)+\sum_{i\ne j}P(c_i,c_j)-h\nu_F.
\end{split}
\end{equation}
The sum over \(i\ne j\) is ordered. Summing these equations gives
\((\D)F=Q(F,F)\). This identity also holds with a common collision cutoff.
When labels include periodic copies, the individual functions are defined
in the lift and are covariant under translations; their sums are periodic.

Constants denoted \(C\), with or without subscripts, may depend on the
fixed construction parameters. They never depend on the approximation
index or the shrinking time horizon unless explicitly stated. The
parameters needed for branching growth will be fixed before the final
small geometric and time parameters. Section~\ref{sec:branching} records
the order of choices.
